\documentclass[11pt]{article}
\usepackage[margin=1.1in]{geometry}
\usepackage{amsmath,amssymb,amsthm}
\usepackage{booktabs}
\usepackage{microtype}
\usepackage[colorlinks=true,linkcolor=blue,citecolor=blue]{hyperref}

\newtheorem{proposition}{Proposition}
\newtheorem{lemma}{Lemma}
\theoremstyle{definition}
\newtheorem{exercise}{Exercise}
\theoremstyle{remark}
\newtheorem{remark}{Remark}

\newcommand{\dd}{\,\mathrm{d}}
\newcommand{\Om}{\Omega}
\newcommand{\pd}[2]{\frac{\partial #1}{\partial #2}}
\newcommand{\R}{\mathbb{R}}
\newcommand{\code}[1]{\texttt{#1}}

\title{Computing $\|u\|_{L^2}$ from the Boundary\\
\large A primer on \texttt{lappy}'s boundary quadrature}
\author{lappy development notes}
\date{}

\begin{document}
\maketitle

\begin{abstract}
\noindent
This is a working introduction to the mathematics behind
\code{lappy.eigfun\_integrals}, aimed at someone who knows vector calculus and a
little numerical analysis and wants to understand --- and modify --- the code as
it stands today. It develops one identity, shows why discretizing it is harder
than it looks, and explains each of the three things the code has to get right
before it can claim an answer to $10^{-13}$. Everything is checkable: every
number quoted can be reproduced by the snippets in the text, and the exercises at
the end are the ones worth doing first.

\medskip\noindent
The companion document \code{corner\_quadrature.pdf} is the research record --- it
derives the corner rule properly and reports what was tried and discarded. Read
this one first.
\end{abstract}

\tableofcontents

\section{The problem, and why the boundary}

We solve $-\Delta u=\lambda u$ in a planar domain $\Om$ with $u=0$ on
$\partial\Om$. A method-of-particular-solutions (MPS) eigensolver gives us
$\lambda$ and a coefficient vector: $u$ is a finite combination of functions that
each solve the PDE \emph{exactly} in $\Om$, chosen so that $u$ is as close to
zero on $\partial\Om$ as the basis allows. So $u$ is not an approximation in the
finite-element sense; it is an exact solution of the PDE with a small nonzero
boundary trace.

Very often we then need the $L^2$ norm,
\begin{equation}
  \|u\|_{L^2(\Om)}^2 = \int_\Om u^2 \dd A,
\end{equation}
because eigenfunctions are only defined up to scale and nearly every downstream
formula wants them normalized:

\begin{itemize}
\item to return orthonormal eigenfunctions from a degenerate cluster;
\item to certify an eigenvalue --- the Moler--Payne bound is a ratio of a
      boundary sup-norm to an interior $L^2$ norm;
\item to evaluate shape derivatives (Hadamard-type formulas), which are boundary
      integrals of normalized eigenfunctions.
\end{itemize}

The obvious route is a cubature rule on the interior: triangulate $\Om$, sum
$w_i u(x_i)^2$. This works, and \code{lappy} still has it
(\code{lappy.cubature}), but it is the expensive and fragile half of any
computation that uses it. A triangulation is a data structure to build, refine
and store; on a curved boundary the polygonal mesh does not even cover $\Om$
exactly; on a spiral domain the mesh takes longer to build than the eigenvalue
took to find.

The alternative is to compute an \emph{area} integral from \emph{boundary} data
only. That is what the Rellich identity does, and it is the subject of this
primer.

\section{The identity}

\subsection{Statement}

Fix any point $x_0\in\R^2$ --- it is a free parameter, and nothing below assumes
anything about it. Write $r:=x-x_0$ and let $n$ be the outward unit normal.

\begin{proposition}[Rellich--Pohozaev, Dirichlet form]\label{prop:rellich}
Let $u$ solve $-\Delta u=\lambda u$ in $\Om\subset\R^2$ with $u=0$ on
$\partial\Om$. Then
\begin{equation}\label{eq:rellich}
  \int_\Om u^2 \dd A
  \;=\; \frac{1}{2\lambda}\int_{\partial\Om} (r\cdot n)
        \left(\pd{u}{n}\right)^{\!2} \dd s .
\end{equation}
\end{proposition}

Everything on the right is boundary data: the geometry ($r\cdot n$) and the
normal derivative of $u$. No interior evaluation appears at all.

\subsection{Proof, in full}

The trick is to multiply the equation by $r\cdot\nabla u$ --- the derivative of
$u$ in the radial direction from $x_0$ --- and integrate. This is the standard
Pohozaev manoeuvre and it is worth doing once by hand.

\begin{proof}
Both sides of $-\Delta u=\lambda u$ are multiplied by $r\cdot\nabla u$ and
integrated over $\Om$. Take the right-hand side first. Since
$u\,(r\cdot\nabla u)=\tfrac12 r\cdot\nabla(u^2)$, the divergence theorem gives
\begin{equation}
  \lambda\!\int_\Om u\,(r\cdot\nabla u)
  = \frac{\lambda}{2}\int_\Om r\cdot\nabla(u^2)
  = \frac{\lambda}{2}\left[\int_{\partial\Om}(r\cdot n)\,u^2
      - \int_\Om u^2\,\nabla\!\cdot\! r\right]
  = -\lambda\!\int_\Om u^2 ,
\end{equation}
using $u=0$ on $\partial\Om$ to kill the boundary term and $\nabla\cdot r=2$ in
the plane. (In $\R^d$ that $2$ would be $d$; the identity's constants are
dimension-dependent, and \code{lappy} is planar throughout.)

For the left-hand side, integrate by parts once:
\begin{equation}\label{eq:byparts}
  \int_\Om \Delta u\,(r\cdot\nabla u)
  = \int_{\partial\Om}\pd{u}{n}(r\cdot\nabla u)
    - \int_\Om \nabla u\cdot\nabla(r\cdot\nabla u).
\end{equation}
For the last integrand, note $\nabla(r\cdot\nabla u)=\nabla u+(r\cdot\nabla)\nabla u$,
so
\begin{equation}
  \nabla u\cdot\nabla(r\cdot\nabla u)
  = |\nabla u|^2 + \tfrac12\, r\cdot\nabla\big(|\nabla u|^2\big),
\end{equation}
and applying the divergence theorem to the second piece,
\begin{equation}
  \int_\Om \nabla u\cdot\nabla(r\cdot\nabla u)
  = \int_\Om|\nabla u|^2
    + \frac12\int_{\partial\Om}(r\cdot n)|\nabla u|^2
    - \int_\Om|\nabla u|^2
  = \frac12\int_{\partial\Om}(r\cdot n)|\nabla u|^2 .
\end{equation}
Now use the boundary condition. Because $u\equiv0$ along $\partial\Om$, its
tangential derivative vanishes there, so the gradient is purely normal:
$\nabla u=(\partial u/\partial n)\,n$ on $\partial\Om$. Hence
$r\cdot\nabla u=(r\cdot n)\,\partial u/\partial n$ and
$|\nabla u|^2=(\partial u/\partial n)^2$, and \eqref{eq:byparts} collapses to
\begin{equation}
  \int_\Om \Delta u\,(r\cdot\nabla u)
  = \int_{\partial\Om}(r\cdot n)\left(\pd{u}{n}\right)^{\!2}
    - \frac12\int_{\partial\Om}(r\cdot n)\left(\pd{u}{n}\right)^{\!2}
  = \frac12\int_{\partial\Om}(r\cdot n)\left(\pd{u}{n}\right)^{\!2}.
\end{equation}
Since $\int(-\Delta u)(r\cdot\nabla u)=\lambda\int u(r\cdot\nabla u)$, the two
computations give
$-\tfrac12\int_{\partial\Om}(r\cdot n)(\partial u/\partial n)^2
 = -\lambda\int_\Om u^2$, which is \eqref{eq:rellich}.
\end{proof}

\subsection{Three consequences worth internalizing}

\paragraph{(i) $x_0$ is free, and that is a gift.} The proof never used any
property of $x_0$. So \eqref{eq:rellich} holds for every choice, and if you
evaluate the right-hand side numerically at two different $x_0$ and get two
different answers, the difference is \emph{error} --- you have a diagnostic that
costs nothing and needs no reference solution. \code{lappy} uses this, with a
caveat we return to in \S\ref{sec:instruments}.

\paragraph{(ii) A corner can be made to disappear.} If $x_0$ sits exactly at a
corner, then on both straight edges meeting there, $r=x-x_0$ points \emph{along}
the edge while $n$ is perpendicular to it, so $r\cdot n\equiv0$ on those edges.
Whatever nasty behaviour $\partial u/\partial n$ has at that corner is multiplied
by an exact zero. This is why \code{default\_x0} places $x_0$ on a singular
corner when there is one. It handles one corner exactly and for free --- but
$x_0$ is a single point, so a domain with several bad corners cannot be rescued
this way, which is the entire reason the rest of this document exists.

\paragraph{(iii) Setting $u\equiv1$ is not allowed, but the same trick is.} You
cannot test \eqref{eq:rellich} with $u=1$ (it is not an eigenfunction). But the
divergence theorem alone gives a closely related identity that is pure geometry,
\begin{equation}\label{eq:area}
  \int_{\partial\Om}(r\cdot n)\dd s = \int_\Om \nabla\cdot r \dd A = 2\,|\Om| ,
\end{equation}
again for \emph{any} $x_0$. Equation \eqref{eq:area} contains no $\lambda$, no
eigenfunction and no basis, so it isolates the quadrature's treatment of the
\emph{geometry}. It turns out to be the sharpest single test in the whole
subject; \S\ref{sec:geom} is about what it revealed.

\subsection{A worked example you should reproduce}

Take the unit disk, $\Om=\{|x|<1\}$, and $x_0=0$, so $r\cdot n=1$ everywhere on
$\partial\Om$. The Dirichlet eigenfunctions are
$u=J_m(kr)\cos(m\theta)$ with $k=j_{m,n}$, the $n$-th positive zero of $J_m$.
On the boundary $\partial u/\partial n = k\,J_m'(k)\cos(m\theta)$, so the right
side of \eqref{eq:rellich} is
\begin{equation}
  \frac{1}{2k^2}\int_0^{2\pi} k^2 J_m'(k)^2\cos^2(m\theta)\dd\theta
  = \frac{\pi}{2}J_m'(j_{m,n})^2
  = \frac{\pi}{2}J_{m+1}(j_{m,n})^2 ,
\end{equation}
using $J_m'(j_{m,n})=-J_{m+1}(j_{m,n})$, and that is exactly the classical
formula for $\|u\|^2$ on the disk. The identity is not merely true, it is
\emph{recognizably} true on the one domain where you know everything.

\section{Discretizing: panels and nodes}

To evaluate \eqref{eq:rellich} we need a quadrature rule on $\partial\Om$: nodes
$x_i$, weights $w_i$, normals $n_i$, and
$\int f \dd s \approx \sum_i w_i f(x_i)$.

\code{lappy} builds this once per solve and reuses it for every $\lambda$,
because the node set depends only on the geometry, on an upper bound
$\lambda_{\max}$ for the eigenvalues of interest, and on a requested accuracy ---
never on the basis or on $u$. That is the design constraint to keep in mind: a
routine that needs to look at $u$ to place nodes cannot live here.

The boundary is a list of \emph{segments} (a line, an arc, a spline), each
parametrized by normalized arclength $\tau\in[0,1]$, so that a node at $\tau$
sits at arclength $\tau\cdot L_{\text{seg}}$ and the weights carry no per-node
Jacobian. Each segment is cut into \emph{panels}, and each panel gets its own
rule:

\begin{center}
\begin{tabular}{ll}
\toprule
\code{legendre} & plain Gauss--Legendre; used on smooth stretches \\
\code{cornerjac}, \code{cornerinterp} & corner-adapted rules, anchored at a corner\\
\bottomrule
\end{tabular}
\end{center}

\noindent
The entry point is
\begin{verbatim}
    from lappy.eigfun_integrals import boundary_quadrature
    bq = boundary_quadrature(domain, lam_max, precision=1e-13)
\end{verbatim}
and \code{bq} carries \code{pts}, \code{normals}, \code{tangents}, \code{wts},
its own reference point \code{x0}, the \code{panels} it was built from, and two
fields that report honesty rather than data: \code{precision} and
\code{shortfalls} (\S\ref{sec:hard}).

Given \code{bq}, evaluating the identity is three lines:
\begin{verbatim}
    from lappy.eigfun_integrals import eigfun_cauchy_data, gram
    ed = eigfun_cauchy_data(basis, lam, coef, bq)   # u, du/dn, du/dt at nodes
    G  = gram(ed, lam, bq)                          # (mult x mult) Gram matrix
\end{verbatim}
For a single eigenfunction \code{G} is $1\times1$ and equals $\|u\|^2$. For a
degenerate cluster of multiplicity $m$ it is the $m\times m$ matrix of inner
products $\langle u_i,u_j\rangle$, from which \code{lowdin\_transform} produces
an orthonormal combination. Note the order of operations: the basis is contracted
with the coefficients \emph{first}, and only then are integrals formed. Never
assemble a basis-level Gram matrix and sandwich it --- \S\ref{sec:scope}.

\section{Why this is not just ``call Gauss--Legendre''}

If $\partial\Om$ is smooth and $u$ is smooth, Gauss--Legendre converges
spectrally and there is no story. The story is corners.

\subsection{What an eigenfunction looks like at a corner}

Put a corner at the origin with interior angle $\alpha$, edges along $\theta=0$
and $\theta=\alpha$, and define
\begin{equation}
  \nu := \frac{\pi}{\alpha}.
\end{equation}
So a right angle has $\nu=2$, a straight-through ``corner'' $\nu=1$, a
$270^\circ$ reentrant corner $\nu=2/3$, and a slit ($\alpha\to2\pi$) $\nu\to1/2$.
Every Dirichlet solution has the local expansion
\begin{equation}\label{eq:expansion}
  u(r,\theta)=\sum_{k\ge1}c_k\,J_{k\nu}\!\big(\sqrt\lambda\,r\big)\sin(k\nu\theta),
\end{equation}
which is worth reading twice: the angular functions $\sin(k\nu\theta)$ are
exactly those vanishing on both edges, and the radial factors are then forced to
be Bessel functions of matching order.

On the edge $\theta=0$ the normal derivative is $-\frac1r\partial_\theta u$, and
using $J_\mu(z)=(z/2)^\mu A_\mu(z^2)$ with $A_\mu$ entire and $A_\mu(0)\ne0$,
\begin{equation}\label{eq:dudn}
  \pd{u}{n}\bigg|_{\theta=0}= r^{\nu-1}F(r),
  \qquad\text{so}\qquad
  \left(\pd{u}{n}\right)^{\!2} = r^{2\nu-2}G(r),
\end{equation}
where $F$ and $G$ are series in powers $r^{k\nu}$ and $r^{2q}$. \textbf{This is
the single most important formula in this document.} The integrand of
\eqref{eq:rellich} near a corner is a smooth series multiplied by $r^{2\nu-2}$.

\subsection{Gauss--Legendre on $r^\gamma$}

So the question is: how well does a Gauss rule integrate $\tau^\gamma$ with
$\gamma=2\nu-2$? Write it down and try it:
\begin{verbatim}
    from lappy.quad import smooth_power_error
    smooth_power_error(gamma, order)   # relative error on int_0^1 tau^gamma
\end{verbatim}

\begin{center}
\begin{tabular}{lccccc}
\toprule
$\gamma=2\nu-2$ & $\nu$ & $\alpha$ & $n=16$ & $n=64$ & fitted rate \\
\midrule
$1/3$    & $7/6$   & $154^\circ$ & $8.0\times10^{-5}$  & $2.1\times10^{-6}$  & $n^{-2.6}$\\
$2/3$    & $4/3$   & $135^\circ$ & $1.2\times10^{-5}$  & $1.2\times10^{-7}$  & $n^{-3.3}$\\
$1$      & $3/2$   & $120^\circ$ & $0$                 & $0$                 & exact\\
$3/2$    & $7/4$   & $103^\circ$ & $1.7\times10^{-7}$  & $1.8\times10^{-10}$ & $n^{-5.0}$\\
$11.55$  & $6.78$  & $27^\circ$  & $8.7\times10^{-16}$ & $3.5\times10^{-16}$ & exact\\
\bottomrule
\end{tabular}
\end{center}

\medskip\noindent
Three things to take from this table.

\paragraph{Integer $\gamma$ is exact.} An $n$-point Gauss rule integrates
polynomials of degree $\le 2n-1$ exactly, and $\tau^\gamma$ with integer $\gamma$
is a polynomial. So $\nu\in\{1,\tfrac32,2,\tfrac52,3,\dots\}$ --- which includes
every right angle and every $120^\circ$ corner --- costs nothing at all.

\paragraph{Non-integer $\gamma$ converges algebraically, at rate $n^{-(2\gamma+2)}$.}
The function $\tau^\gamma$ is continuous but not analytic at $0$: some derivative
blows up. Gauss--Legendre's spectral convergence needs analyticity, and without
it the classical result is an algebraic rate governed by how many derivatives
exist. The fitted rates above match $2\gamma+2$ to two digits ($2.64$ vs $2.67$,
$3.31$ vs $3.33$, $4.96$ vs $5.00$), which is a satisfying check to run yourself.

\paragraph{The rate is what matters, not the exponent's sign.} At
$\gamma=11.55$ the predicted rate is $n^{-25}$: the function is so smooth that
Gauss is at roundoff by order 16, even though $\gamma$ is irrational. At
$\gamma=2/3$ the rate is $n^{-3.3}$ and you will \emph{never} reach $10^{-13}$
--- order 64 gives $10^{-7}$, and quadrupling again buys three digits.

\subsection{The criterion \texttt{lappy} uses, and the one it used to use}

Here is the practical payoff, and a mistake worth learning from. For a long time
the code decided a corner ``needs the special rule'' when $\nu<1$, i.e.\ when the
corner is \emph{reentrant}, on the reasoning that a convex corner has bounded $u$
and bounded $\partial u/\partial n$, so a smooth rule should be fine.

That reasoning is wrong, and the table above says why. At a $135^\circ$ corner,
$\nu=4/3$: the eigenfunction is bounded, its normal derivative $r^{1/3}$ even
\emph{vanishes} at the corner --- and the integrand is $r^{2/3}$, which Gauss
integrates at rate $n^{-3.3}$. Nothing is singular in the everyday sense, and the
quadrature still fails. What matters is not whether $2\nu-2$ is negative but
whether it is an \emph{integer}.

Equally, the opposite over-correction is also wrong: ``use the corner rule
whenever $\nu$ is not an integer'' hands the special rule to $\nu=6.78$ corners
that do not need it and makes them \emph{worse}, because the corner rule is built
to be exact on a singular class and is not even exact on constants
(\S\ref{sec:looks-like-bugs}). So \code{lappy} asks the question directly ---
what does a smooth rule achieve on $\tau^{2\nu-2}$ at the order this segment is
getting? --- and uses the corner rule exactly when the answer is not good enough.
Measured on real domains, by refining the node set and differencing the integral:

\begin{center}
\begin{tabular}{lccc}
\toprule
domain & nodes & before & after\\
\midrule
\code{right\_trapezoid} & $84\to212$ & $1.8\times10^{-6}$ & $2.9\times10^{-16}$\\
\code{GWW1} & $204\to440$ & $3.2\times10^{-5}$ & $5.0\times10^{-13}$\\
\code{iso\_tri\_h05} & $96\to224$ & $1.2\times10^{-11}$ & $4.8\times10^{-16}$\\
\bottomrule
\end{tabular}
\end{center}

\section{The corner-adapted rule, in one page}

When a corner does need help, the fix is to build a rule matched to the class
\eqref{eq:dudn} rather than to polynomials. Two ideas, both standard:

\paragraph{Absorb the singularity into the weight.} Gauss--Jacobi rules integrate
$\int_0^1\tau^{\gamma}(\text{polynomial})\dd\tau$ exactly, with $\gamma$ as a
weight exponent. Taking $\gamma=2\nu-2$ handles the $r^{2\nu-2}$ prefactor by
construction, leaving a smooth series for the polynomial part.

\paragraph{Rationalize the exponents.} $G$ in \eqref{eq:dudn} contains powers
$r^{k\nu}$ \emph{and} $r^{2q}$ --- two incommensurate families. A substitution
$\tau=t^{q}$ with $\nu=p/q$ turns both into integer powers of $t$, so a
polynomial rule in $t$ is exact on the whole class. This works beautifully in
theory and fails in practice for $q\ge4$: the innermost node lands at
$t_{\min}^q$, which underflows the coordinate representation (a node at
$p_0+\tau L\hat d$ rounds onto $p_0$ itself), and the accuracy is destroyed.
\code{lappy} therefore takes the node \emph{placement} from the mild exponent
$\nu$ and gets exactness from the \emph{weights} instead, via an interpolatory
rule on the true exponent set. That is the difference between \code{cornerjac}
(substitution; only safe for straight edges with $2/\nu$ integral) and
\code{cornerinterp} (the general case).

Two practical consequences follow, and both look like bugs the first time.

\subsection{Two things that look like bugs}\label{sec:looks-like-bugs}

\paragraph{The weights do not sum to the panel length.} A corner rule is exact on
$\{\gamma+k\nu+m\}$, a set that does not contain $\tau^0$. It therefore does not
integrate constants, and $\sum_i w_i \ne L$ on that panel. Renormalizing would
destroy exactness on the class the rule exists for. (This is also why feeding a
corner rule an integrand with no corner structure --- a plane wave, say --- gives
a bad answer, and why a naive ``always use the corner rule'' criterion loses.)

\paragraph{Accuracy is not monotone in the order.} Past a $\nu$-dependent
threshold, raising the order makes things \emph{worse}: the integrand's dynamic
range ($\tau^{2\nu-2}\sim6\times10^8$ at $\tau\sim10^{-9}$) against
correspondingly tiny weights means roundoff, not truncation, dominates. Order
selection therefore \emph{scans} a calibrated curve for the best order; it never
increases the order until a target is met.

\section{Three requirements, not one}\label{sec:sizing}

Choosing each panel's order is where the code spends its cleverness. A panel must
resolve all of the following, and the required order is the largest of the three.

\paragraph{1. Oscillation.} An eigenfunction oscillates at wavenumber
$\sqrt\lambda$, and the integrand is a \emph{product} of two derivatives, so it
oscillates at $2\sqrt\lambda$. Over a segment of length $L$ that is a total phase
$k=2\sqrt{\lambda_{\max}}\,L$, and \code{smooth\_order\_for\_precision} finds the
smallest Gauss order integrating $e^{ik\tau}$ to the target. Model integrand,
closed-form answer, so the sizing is self-certifying rather than calibrated
offline.

\paragraph{2. The corner class.} As \S4: if the smooth rule cannot handle
$\tau^{2\nu-2}$, the panel is corner-anchored and its order comes from a model
integrand in the corner's own exponent family
(\code{corner\_order\_for\_precision}).

\paragraph{3. The parametrization.}\label{sec:geom} This one is easy to forget,
and it is the subject of the newest part of the code. Panels live in
\emph{normalized arclength}. For a straight edge or a circular arc, that
parametrization is exact and free. For a curve of varying speed --- an ellipse
--- the arclength map is analytic but its strip of analyticity narrows as the
curve becomes eccentric, and Gauss--Legendre's convergence rate narrows with it.

You can see this with no eigenfunction at all, using the pure-geometry identity
\eqref{eq:area}:

\begin{center}
\begin{tabular}{lcccccc}
\toprule
order & 16 & 32 & 64 & 128 & 192 & 256\\
\midrule
disk          & $3$e$-16$ & $3$e$-16$ & $3$e$-16$ & $3$e$-16$ & $3$e$-16$ & $3$e$-16$\\
ellipse $a=2$ & $5$e$-03$ & $2$e$-04$ & $4$e$-07$ & $3$e$-12$ & $5$e$-15$ & $2$e$-15$\\
ellipse $a=4$ & $4$e$-02$ & $1$e$-02$ & $2$e$-03$ & $8$e$-05$ & $3$e$-06$ & $1$e$-07$\\
ellipse $a=6$ & $6$e$-02$ & $3$e$-02$ & $1$e$-02$ & $2$e$-03$ & $5$e$-04$ & $1$e$-04$\\
\bottomrule
\end{tabular}
\end{center}

\medskip\noindent
A circle is exact at every order because its arclength map is linear. An
$a=2$ ellipse needs about $170$ nodes where the oscillation model asked for $46$.
An $a=6$ ellipse cannot be resolved at any order you would want to use.
\code{geometry\_order\_for\_precision} sizes each segment on
\eqref{eq:area}, so this is caught automatically --- and costs nothing on
every exactly-parametrized boundary, because there the geometry order never
exceeds the oscillation order.

\begin{remark}
A natural guess is that the culprit is the interpolation tolerance of the
arclength table. It is not: the error is identical to three digits at
$\text{tol}=10^{-4}$, $10^{-6}$ and $10^{-8}$, while the table's build cost goes
$0.0$s, $0.9$s, $65$s. Refining it is pure expense. The lesson is worth
generalizing --- before optimizing a suspected cause, vary it and check that the
symptom responds.
\end{remark}

\section{What is out of reach, and how the code says so}\label{sec:hard}

Some geometry is genuinely hard, and the design goal is not to pretend otherwise.
Three examples, all in the benchmark suite:

\begin{itemize}
\item a \textbf{near-slit corner} ($\nu\to\tfrac12$), where the integrand
      $r^{2\nu-2}\to r^{-1}$ is barely integrable and $90\%$ of the integral
      lives within $10^{-9}$ of the corner;
\item an \textbf{eccentric ellipse}, per \S\ref{sec:sizing};
\item a \textbf{coiled spiral}, with many corners and long thin channels.
\end{itemize}

For these, \code{boundary\_quadrature} does three things instead of quietly
returning a wrong number:
\begin{enumerate}
\item \code{bq.precision} reports what the rule can support, not what was asked
      for. A singular corner that had to be demoted to a smooth rule gets
      $\infty$, because there is no bound at all in that case.
\item \code{bq.shortfalls} names each corner or segment that fell short, and why.
\item a warning distinguishes the \emph{cause} --- in particular an eccentric
      ellipse warns that the \emph{parametrization} is the limit, so you do not
      go hunting for a bug in the eigenfunction.
\end{enumerate}

There is a fourth class, currently labelled rather than solved: a stadium has
exact parametrization and integer $\nu$, yet its integrals sit near
$10^{-4}$, because its boundary data carries far more harmonics than
$2\sqrt\lambda L$ predicts. No geometry-only rule can anticipate that --- it is a
property of the solution, not of the domain --- so the honest tool there is the
a posteriori check of the next section.

\section{How to check your own work}\label{sec:instruments}

\paragraph{Refine and difference (the reliable one).}
\begin{verbatim}
    from lappy.eigfun_integrals import verify_gram
    G, G_ref, err = verify_gram(basis, lam, coef, bq, domain)
\end{verbatim}
recomputes the Gram on a systematically refined node set and returns the change.
This measures the quadrature and nothing else. It is the instrument to reach for.

\paragraph{Vary $x_0$ (the free detector, with a caveat).} By
Proposition~\ref{prop:rellich} the answer cannot depend on $x_0$, so evaluating
\code{gram(ed, lam, bq, x0=...)} at several points and looking at the spread
costs three matrix products and needs no reference. It is what first exposed the
$\nu<1$ mistake. But it is a detector, not an error bar: it also carries the
\emph{eigenvalue's} error, since $\lambda$ enters through $1/(2\lambda)$ and
through the Cauchy data, so an inexact $\lambda$ makes $u$ not quite an
eigenfunction and the $x_0$-dependence reappears. Measured, the two can differ by
three orders in either direction --- on one domain the spread reads
$3.3\times10^{-12}$ where the quadrature error is $1.4\times10^{-13}$; on another
it reads $4.1\times10^{-11}$ where the quadrature error is $1.0\times10^{-8}$.

\paragraph{Compare against closed forms (the ground truth).} \code{lappy.reference}
has exact eigenfunctions and exact $L^2$ norms for rectangles, isosceles right
triangles, disks, circular sectors of arbitrary angle, and polyominoes. These are
the only unimpeachable references, and every test in
\code{tests/test\_eigfun\_integrals.py} uses one.

\begin{remark}[on instruments]
Do not use \code{mpmath.quad} as a reference for these integrands. As
$\nu\to\tfrac12$ the endpoint behaviour approaches $r^{-1}$ and it errs by
$4\times10^{-2}$ even at 40 digits --- manufacturing convincing, entirely
spurious, convergence plateaus. This cost a real amount of time.
\end{remark}

\section{Scope: what this machinery is \emph{not} for}\label{sec:scope}

The rule is matched to \emph{eigenfunction}-level data. At a corner, an exact
eigenfunction's local structure is purely the corner family \eqref{eq:expansion},
which is exactly what the rule integrates.

It is not appropriate for a basis-level $N\times N$ Gram matrix. A basis column
centred at some other corner is plain analytic at this one, with $O(1)$ amplitude
at every corner simultaneously --- the wrong class entirely. Measured, the corner
rule beats a graded rule by 3--8 orders on corner-family blocks and \emph{loses}
by 2--4 on mixed ones, which is why \code{lappy} retired the basis-level path
rather than adapting it. Hence the slogan in the module docstring: evaluate
first, sandwich never. Contract the basis with the coefficients, \emph{then}
integrate.

\section{Exercises}

These are in the order I would do them. Each should take minutes, and each has a
number you can check against.

\begin{exercise}[the geometry identity]
For \code{disk(1.0)}, \code{rect(1,1)} and \code{ellipse(2,1)}, build a node set
and evaluate $\sum_i w_i\,(x_i-x_0)\cdot n_i$ against $2|\Om|$, for two or three
different $x_0$. Which domains are exact, and which is not? Then re-run the
ellipse with \code{resolve\_geometry=False} and explain the difference.
\end{exercise}

\begin{exercise}[the disk, in closed form]
Using \code{lappy.reference.disk\_eigfun}, form the Cauchy data of $J_2(j_{2,1}r)\cos2\theta$
on the boundary by hand, evaluate \eqref{eq:rellich} with the quadrature, and
compare to $\tfrac{\pi}{2}J_3(j_{2,1})^2$. You should get $10^{-15}$ or better.
Now do it for a circular sector of angle $\tfrac32\pi$ with
\code{sector\_eigfun}, whose apex is a reentrant corner, and try $x_0$ at the
apex versus $x_0$ elsewhere.
\end{exercise}

\begin{exercise}[the rate law]
Reproduce the $\tau^\gamma$ table with \code{smooth\_power\_error} and fit the
convergence rate for $\gamma=1/3,\,2/3,\,3/2$. Confirm $n^{-(2\gamma+2)}$. Then
try $\gamma=4.5$ and explain why the fit breaks down. (Hint: look at the size of
the numbers, not the trend.)
\end{exercise}

\begin{exercise}[the criterion]
Build a regular octagon ($\nu=4/3$ at every corner) and compare
\code{boundary\_quadrature(..., nonintegral=False)} with the default. Inspect
\code{bq.panels} to see which rules were assigned, then use \code{verify\_gram}
on a real MPS eigenfunction to measure what changed. Predict the answer for a
regular hexagon ($\nu=3/2$) \emph{before} running it.
\end{exercise}

\begin{exercise}[reading a refusal]
Build a node set for \code{ellipse(6,1)} and for a sector of angle $1.984\pi$.
Read \code{bq.precision}, \code{bq.shortfalls} and the warnings. For each, say in
one sentence which of the three requirements in \S\ref{sec:sizing} failed, and
whether more nodes would help.
\end{exercise}

\begin{exercise}[the whole pipeline]
Take \code{L\_shape()}, build an \code{MPSEigensolver} with
\code{from\_domain}, solve for the first eigenvalue, and check that
\code{eigenfunction\_coef} returns coefficients whose Gram is $1$ to
$10^{-13}$. Then verify the same norm against interior cubature
(\code{lappy.cubature.polygon\_cubature}) and explain which of the two you would
trust if they disagreed at $10^{-9}$.
\end{exercise}

\section{Answers to check against}

Run on the code as of this writing. If you get these, you have wired things up
correctly; if you get something else, the discrepancy is the interesting part.

\medskip
\begin{center}
\begin{tabular}{lll}
\toprule
Exercise & quantity & value \\
\midrule
1 & disk, rect: $|\!\sum w\,(r\cdot n)-2|\Om||$ & $\le3\times10^{-15}$, any $x_0$\\
1 & ellipse $(2,1)$, default & $1.4\times10^{-15}$ at 168 nodes\\
1 & ellipse $(2,1)$, \code{resolve\_geometry=False} & $1.5\times10^{-5}$ at 46 nodes\\
\midrule
2 & disk $J_2$ mode, quadrature & $0.181230398151804$\\
2 & $\tfrac{\pi}{2}J_3(j_{2,1})^2$ & $0.181230398151804$\\
\midrule
3 & fitted rates, $\gamma=\tfrac13,\tfrac23,\tfrac32$ & $2.64,\;3.31,\;4.96$\\
3 & predicted $2\gamma+2$ & $2.67,\;3.33,\;5.00$\\
\midrule
4 & octagon ($\nu=4/3$), default & \code{cornerinterp}, 824 nodes\\
4 & octagon, \code{nonintegral=False} & \code{legendre}, 96 nodes\\
4 & hexagon ($\nu=3/2$), either way & \code{legendre}, 84 nodes\\
\midrule
5 & \code{ellipse(6,1)} & \code{precision} $=9.9\times10^{-5}$, segment shortfall\\
5 & sector of $1.984\pi$ & \code{precision} $=\infty$, corner demoted\\
\midrule
6 & \code{L\_shape}, boundary Gram & $|G-1| = 2.2\times10^{-16}$\\
6 & \code{L\_shape}, interior cubature & $|{\cdot}-1| = 5.9\times10^{-11}$\\
\bottomrule
\end{tabular}
\end{center}

\medskip\noindent
Exercise 3's fit breaks down at $\gamma=4.5$ because the error has already reached
roundoff ($\sim10^{-15}$) by order 32; there is no trend left to fit. Exercise 4's
hexagon is the control: $\nu=3/2$ gives $\gamma=1$, an exact polynomial, so
nothing should change and nothing does. Exercise 5's two failures are deliberately
different --- one names a \emph{segment} and blames the parametrization, the other
names a \emph{corner} and refuses to quote a number at all. Exercise 6 is the one
to think hardest about: the boundary route is five orders better than the interior
cubature it is being compared to, which is why the comparison is reported as a
cross-check and not as an error bar.

\section*{Where to read next}

\begin{itemize}
\item \code{docs/corner\_quadrature.pdf} --- the research record: the corner rule
      derived properly, the failure of the graded rule quantified, and an
      addendum listing what did not survive contact with real domains.
\item \code{docs/rellich.md} --- the identity in its general (Zaremba) form,
      covering mixed boundary conditions.
\item \code{docs/eigfun\_integrals.md} --- the API and its accuracy record.
\item \code{lappy/eigfun\_integrals.py} --- the implementation. The docstrings
      carry the measurements; read \code{boundary\_quadrature} and
      \code{corner\_specs} first.
\item \code{tests/test\_eigfun\_integrals.py} --- worked, checked examples of
      every claim above, against closed forms.
\end{itemize}

\end{document}
